
\section{Related Work}

\subsection{Low-Rank Matrix Factorization}

\textbf{Classical SVD.} The Eckart--Young theorem (Eckart \& Young, 1936) guarantees 
that the best rank-$r$ approximation of $W$ in Frobenius norm is the truncated 
SVD:

\[
W_r^* = U_{:, :r} \Sigma_{:, :r} V_{:r, :}^\top
\tag{2.1}
\]

This is optimal in an unweighted sense and forms the basis of most 
SVD-based compression methods.

\textbf{Randomized SVD.} Halko et al. (2011) provide a randomized algorithm for 
truncated SVD with complexity $\mathcal{O}(d_\text{in} d_\text{out} \log r)$, 
which we do not use because exact SVD is feasible for the layer sizes in our 
study.

\textbf{Weighted low-rank approximation.} Golub et al. (1987) and Srebro \& Jaakkola 
(2003) study weighted variants of low-rank approximation. Our Fisher weighting 
is a diagonal special case with input-dimension-dependent weights.

\subsection{Transformer Compression}

\subsubsection{Post-Training Pruning}

Magnitude pruning (Han et al., 2015) removes individual weights. Structured 
pruning (Li et al., 2017) removes entire channels. SparseGPT (Frantar \& 
Alistarh, 2023) and Wanda (Sun et al., 2024) prune large language models 
without retraining. A3 differs by factorizing rather than zeroing.

\subsubsection{Knowledge Distillation}

DistilBERT (Sanh et al., 2019), TinyBERT (Jiao et al., 2020), and MiniLM 
(Wang et al., 2020) train smaller models to mimic larger ones. A3 is 
orthogonal -- it can compress an already-distilled model.

\subsubsection{Quantization}

8-bit (Dettmers et al., 2022) and 4-bit (Dettmers et al., 2023) quantization 
reduces bit-width. GPTQ (Frantar et al., 2023) uses second-order 
information. A3 is orthogonal to quantization -- the two can be combined.

\subsubsection{Low-Rank Factorization}

ALBERT (Lan et al., 2020) uses cross-layer parameter sharing, a form of 
low-rank structure. Tensor train decomposition (Novikov et al., 2015) uses 
higher-order factorization. **A3 differs by using activation-aware weighting 
during factorization**, not after.

\subsection{Activation-Aware Compression}

\textbf{AWQ.} Lin et al. (2023) find that activation magnitude correlates with 
weight importance. They use it for quantization. Our Fisher weighting uses the 
same insight for low-rank factorization.

\textbf{Activation-aware pruning.} Molchanov et al. (2019) use the Taylor expansion 
of the loss with respect to activations. Our approach is simpler: we use 
second moments.

\textbf{Fisher information.} The Fisher diagonal is the classical measure of 
parameter importance (Pascanu \& Bengio, 2014; Kirkpatrick et al., 2017). We 
use a cheap empirical proxy $\mathbb{E}[x^2]$ instead of the full Fisher.

\subsection{LayerNorm Recalibration}

\textbf{Post-training recalibration.} Nagel et al. (2019) recalibrate activation 
distributions for quantization. We adapt this idea to post-factorization 
distributional shift.

\textbf{Gamma rescaling.} Gholami et al. (2021) rescale LayerNorm scales after 
quantization. Our $W^* = U \Sigma V^T$ is the analogous operation 
for low-rank factorization.

\subsection{Position of Our Work}

\begin{table}[H]
\centering
\small
\begin{tabular}{lll}
\toprule
\textbf{Axis} & \textbf{Prior work} & \textbf{A3} \\
\midrule
\textbf{Factorization weighting} & Uniform (unweighted SVD) & Fisher diagonal \\
\textbf{Rank selection} & Fixed & Adaptive escalation \\
\textbf{Calibration memory} & Full activation buffer & Streaming O($d^2$) \\
\textbf{Post-factorization fix} & None & LayerNorm recalibration \\
\textbf{Error validation} & Silent & Strict Frobenius bound \\
\textbf{Fallback} & Never & Automatic per-layer \\
\bottomrule
\end{tabular}
\end{table}

\textbf{Key distinction.} Prior low-rank compression methods either (a) select 
rank globally before factorizing, or (b) factorize unweighted. A3 combines 
both improvements: it \textbf{selects rank per-layer based on validation}, and it 
\textbf{weights the factorization by activation statistics}. The result is 
strictly smaller reconstruction error at the same rank, or equivalently, 
smaller rank at the same error.



\subsection{Comparison Table}
